\chapter{Class Diagram}
\label{ch:class}

\section{Purpose}

The class diagram models the static structure of the data the tool operates on:
the three representations a translation passes through, and the entity types
that make up a GENESYS project.

\umlfigure{class-diagram}{Class diagram}{fig:class}

\section{Representations}

A translation is a chain of three representations, each of which knows how to
translate to the next. Reading left to right gives the import direction; reading
right to left gives the export direction.

\begin{longtable}{L{0.24\textwidth}L{0.68\textwidth}}
	\caption{Translation representations}\label{tab:representations}\\
	\toprule
	\textbf{Class} & \textbf{Responsibility} \\
	\midrule
	\endfirsthead
	\toprule
	\textbf{Class} & \textbf{Responsibility} \\
	\midrule
	\endhead

	\texttt{C\#} & The source-code representation. Holds the program as written,
	and translates to and from the abstract syntax. \\
	\addlinespace

	\texttt{Absyn} & The abstract syntax, and the pivot of the design. It belongs
	to neither the language nor the model, so neither end depends on the other
	directly, and a second language becomes a second translation into the same
	abstract syntax rather than a change to anything downstream. \\
	\addlinespace

	\texttt{Genesys Project} & The model representation. Aggregates the entity
	collections \texttt{functions[]}, \texttt{items[]} and \texttt{state[]}, and
	translates to and from the abstract syntax. \\
	\bottomrule
\end{longtable}

Each class exposes a \texttt{translate()} operation. Translation is therefore a
responsibility of the representation rather than of a separate converter, and
the chain is traversed by calling \texttt{translate()} along it.

\section{Project Entities}

A GENESYS project is made up of entities. The three the tool reads and writes
share a common shape: each carries its relationships to other entities and the
targets and attributes attached to those relationships, which is how the GENESYS
schema itself is organised.

\begin{longtable}{L{0.18\textwidth}L{0.30\textwidth}L{0.44\textwidth}}
	\caption{Project entity types}\label{tab:entities}\\
	\toprule
	\textbf{Class} & \textbf{Attributes} & \textbf{Role} \\
	\midrule
	\endfirsthead
	\toprule
	\textbf{Class} & \textbf{Attributes} & \textbf{Role} \\
	\midrule
	\endhead

	\texttt{Function} & \texttt{relationships[]}, \texttt{targetsAndAttributes[]},
	\texttt{parameters[]}
	& A behaviour in the model. Carries parameters in addition to the common
	shape, because a function's inputs and outputs become a generated method's
	signature. \\
	\addlinespace

	\texttt{Item} & \texttt{relationships[]}, \texttt{targetsAndAttributes[]}
	& What flows between functions. \\
	\addlinespace

	\texttt{State} & \texttt{relationships[]}, \texttt{targetsAndAttributes[]}
	& A state in a state machine. Transitions are reached through
	\texttt{relationships[]} rather than being a class of their own. \\
	\bottomrule
\end{longtable}

\section{Design Decisions}

\paragraph{An abstract syntax as the pivot.}
Placing \texttt{Absyn} between the language and the model means neither end
references the other. This is the structural counterpart of the process boundary
in Figure~\ref{fig:block}, where the extraction half and the model half exchange
a neutral representation rather than object references.

\paragraph{Translation belongs to the representation.}
Each representation implements \texttt{translate()} rather than translation
being owned by an external converter. Adding a representation therefore means
adding a class that can translate into the abstract syntax, with no change to
the existing ones.

\paragraph{Entities mirror the GENESYS schema.}
\texttt{Function}, \texttt{Item} and \texttt{State} carry relationships and
targets with attributes because that is how GENESYS stores them. Keeping the
same shape means reading and writing the repository is a direct mapping rather
than a reinterpretation, and that a schema addition does not require a new
class.

\paragraph{Parameters only on Function.}
\texttt{Function} is the only entity carrying \texttt{parameters[]}, since it is
the only one that becomes a method signature. \texttt{Item} and \texttt{State}
reach everything they need through relationships.
